\documentclass[11pt]{article}
\usepackage[margin=1in]{geometry}
\usepackage{amsmath,amssymb,amsthm}
\usepackage[T1]{fontenc}
\usepackage{lmodern}
\usepackage{microtype}
\usepackage[hidelinks]{hyperref}
\newtheorem{theorem}{Theorem}
\newtheorem{corollary}[theorem]{Corollary}
\newcommand{\RR}{\mathbb R}
\newcommand{\tr}{\operatorname{tr}}
\newcommand{\diag}{\operatorname{diag}}
\DeclareMathOperator*{\argmin}{arg\,min}
\title{Resolving User3's orthogonal absolute-product question:\\minimum and minimizers}
\author{Haobo Ma\\\small ChronoAI Pte Ltd, Singapore
\and Wenlin Zhang\\\small National University of Singapore\\\small ChronoAI Pte Ltd, Singapore}
\date{November 1, 2026}
\begin{document}
\maketitle
\begin{abstract}
We resolve User3's question \cite[question 515549, revision 1,
September 25, 2026]{question}: for every nonsingular real square matrix
of positive size, the orthogonal minimum of
$\rho(|PQ|\,|Q^\top P^{-1}|)$ is attained and equals one.
Thus the proposed strict uniform lower bound cannot hold for any such
matrix. We give a self-contained exposition of the triangularization
proof in \cite{formal}, then characterize every minimizer by independent
row and column permutations to triangular form. For fixed $P$, the set
of minimizers is finite and has an explicit description using row-permuted
QR factorizations. The same minimum holds over the special orthogonal
group. Finally, the companion absolute-row-sum balancing question
\cite{balance} follows from the Yamabe--Yujobo theorem.
\end{abstract}

\section{The question and its attained minimum}

For $K\geq1$, write $O(K)=\{Q\in\RR^{K\times K}:Q^\top Q=QQ^\top=I\}$.
For a matrix $A$, $|A|$ denotes entrywise absolute value, and $\rho(A)$
is the maximum modulus of its complex eigenvalues. Fix
$P\in GL_K(\RR)$ and set
\[
 M_P(Q)=|PQ|\,|Q^\top P^{-1}|,\qquad Q\in O(K).
\]
User3 \cite{question} asks whether a quantity $\beta(P)>1$ can bound
$\rho(M_P(Q))$ below for every orthogonal $Q$, on a suitable class of
nonsingular matrices, and asks for an intrinsic characterization of the
infimum. The answer to the second question is the constant one; the answer
to the first is negative for every individual $P$.

The triangularization argument already appears in the formal source
\cite[\texttt{result}]{formal}. We present its complete elementary proof,
using a trace bound in place of spectral-radius monotonicity. Row
Gram--Schmidt is equivalently QR factorization of $P^\top$, rather than
of $P$; see \cite{qr} for the QR construction. Section 2 gives the equality
cases. Section 3 relates the objective to diagonal scaling, and Section 4
settles the companion question from the same posing author.

\begin{theorem}\label{thm:minimum}
For every $K\geq1$ and $P\in GL_K(\RR)$,
\[
 \min_{Q\in O(K)}\rho(M_P(Q))
 =\inf_{Q\in O(K)}\rho(M_P(Q))=1.
\]
An attaining matrix $Q$ is obtained by applying ordered Gram--Schmidt to
the rows of $P$ and placing the resulting unit vectors in its columns.
The same attained minimum holds over $SO(K)$.
\end{theorem}

\begin{proof}
For arbitrary $Q\in O(K)$, put $A=PQ$ and $B=Q^\top P^{-1}$.
Then $AB=I$. The triangle inequality gives, for every $i,j$,
\[
 \delta_{ij}=\left|\sum_{k=1}^K a_{ik}b_{kj}\right|
 \leq\sum_{k=1}^K|a_{ik}|\,|b_{kj}|=(M_P(Q))_{ij}.
\]
In particular, $\tr M_P(Q)\geq K$. If $\lambda_1,\ldots,\lambda_K$
are the eigenvalues, counted with algebraic multiplicity, then
\[
 K\leq\tr M_P(Q)
 =\left|\sum_{j=1}^K\lambda_j\right|
 \leq\sum_{j=1}^K|\lambda_j|
 \leq K\rho(M_P(Q)).
\]
Division by $K>0$ proves the universal bound, without a
diagonalizability assumption.

Let $p_1,\ldots,p_K$ be the row vectors of $P$, viewed in $\RR^K$.
They are linearly independent. Ordered Gram--Schmidt defines
\[
 v_i=p_i-\sum_{j<i}\langle p_i,q_j\rangle q_j,
 \qquad q_i=\frac{v_i}{\|v_i\|_2}.
\]
Each $v_i$ is nonzero by independence. The vectors $q_i$ are an
orthonormal basis and
$p_i\in\operatorname{span}(q_1,\ldots,q_i)$.
For $Q_*=[q_1\ \cdots\ q_K]$, the matrix $T=PQ_*$ therefore satisfies
\[
 t_{ij}=\langle p_i,q_j\rangle=0\quad(i<j),
 \qquad t_{ii}=\|v_i\|_2>0.
\]
Thus $T$ is nonsingular lower triangular, and
$T^{-1}=Q_*^\top P^{-1}$ is lower triangular. For completeness,
forward substitution in $Tx=e_j$ gives $x_i=0$ for $i<j$ and
$x_j=1/t_{jj}$, establishing both assertions about the inverse.

The matrix $|T|\,|T^{-1}|$ is lower triangular. In its $i$th diagonal
entry, the first factor can be nonzero only for $k\leq i$, while the
second can be nonzero only for $k\geq i$. Hence
\[
 (|T|\,|T^{-1}|)_{ii}
 =|t_{ii}|\,|(T^{-1})_{ii}|=1.
\]
Its characteristic polynomial is $(x-1)^K$, so its spectral radius is
one. This proves actual attainment, and therefore the minimum and
infimum assertions.

Finally, for any diagonal sign matrix $D$, entrywise absolute values give
$M_P(QD)=M_P(Q)$. Choose $D$ with $\det D=\det Q_*$; such a choice is
possible also when $K=1$. Then $Q_*D\in SO(K)$ and still attains one.
\end{proof}

Consequently there is no $\beta(P)>1$ valid for every orthogonal $Q$, for any
nonsingular $P$ or any nonempty class of such matrices. This settles both
requests in \cite{question}, including attainment rather than just an
asymptotic approximation to the infimum.

\section{Every equality case and every minimizer}

The distinction between a unit spectral radius and the identity matrix
is essential. The following characterization makes that distinction
precise. Row and column permutations in its statement are independent.

\begin{theorem}\label{thm:equality}
For $A\in GL_K(\RR)$, $K\geq1$,
\[
 \rho(|A|\,|A^{-1}|)=1
 \quad\Longleftrightarrow\quad
 SAC\text{ is lower triangular for some permutation matrices }S,C.
\]
\end{theorem}

\begin{proof}
Put $M=|A|\,|A^{-1}|$. As in the preceding proof, $M\geq I$ entrywise.
Suppose $\rho(M)=1$. Every eigenvalue $\lambda_j$ has
$\Re\lambda_j\leq1$, whereas
\[
 K\leq\tr M=\sum_{j=1}^K\Re\lambda_j\leq K.
\]
Equality forces every $\Re\lambda_j=1$; together with $|\lambda_j|\leq1$
this gives $\lambda_j=1$. Thus $N=M-I\geq0$ has characteristic polynomial
$x^K$, and Cayley--Hamilton gives $N^K=0$.

The directed graph with an edge $i\to j$ when $N_{ij}>0$ has no directed
cycle. A positive cycle of length $r$ would give
$(N^{mr})_{ii}>0$ for every positive integer $m$, contradicting
nilpotence. Reverse topological ordering therefore gives a permutation
$S$ for which $M'=SMS^\top$ is lower triangular with diagonal one.
Set $A'=SA$ and $B'=A^{-1}S^\top$, so $M'=|A'|\,|B'|$.
For each $j$, choose $k_j$ with
\[
 |A'_{jk_j}|\,|B'_{k_jj}|>0;
\]
such an index exists because $M'_{jj}=1$.
These indices are distinct: if $i<j$ and $k_i=k_j$, then
$|A'_{ik_i}|\,|B'_{k_jj}|>0$ contributes to the zero entry $M'_{ij}$,
a contradiction. Furthermore, for $i<j$, nonnegativity and
$M'_{ij}=0$ imply $A'_{ik_j}=0$, since $B'_{k_jj}\ne0$.
The permutation matrix $C=[e_{k_1}\ \cdots\ e_{k_K}]$ consequently
makes $SAC$ lower triangular.

Conversely, if $T=SAC$ is nonsingular lower triangular, then
\[
 |T|\,|T^{-1}|=S|A|\,|A^{-1}|S^\top
\]
is lower triangular with diagonal one, as in Theorem~\ref{thm:minimum}.
Permutation similarity preserves the spectrum, proving the converse.
\end{proof}

\begin{corollary}\label{cor:minimizers}
For each row permutation $S$, fix the factorization
$SPU_S=L_S$ with $U_S$ orthogonal and $L_S$ lower triangular with
positive diagonal. Then
\[
 \argmin_{Q\in O(K)}\rho(M_P(Q))
 =\{U_SW:S\text{ a permutation matrix},\ W\text{ a signed permutation matrix}\}.
\]
In particular, there are at most $2^K(K!)^2$ minimizing matrices.
\end{corollary}

\begin{proof}
A signed permutation matrix has exactly one entry $\pm1$ in each row
and column. Its absolute value is a permutation matrix, so
$M_P(QW)=M_P(Q)$.
If $Q$ minimizes, Theorem~\ref{thm:equality} gives $S,C$ with
$T=SPQC$ lower triangular. Then
\[
 U_S^\top QC=L_S^{-1}T
\]
is both orthogonal and lower triangular. Such a matrix is diagonal with
diagonal entries $\pm1$: its last column has only its last entry
nonzero, with modulus one; orthogonality makes the other entries in the
last row zero, and induction completes the assertion.
Thus $Q=U_SDC^\top$ for a diagonal sign matrix $D$, which has the required
form. Conversely, $SPU_S=L_S$ and permutation similarity give
$\rho(M_P(U_S))=1$; signed-permutation invariance gives the same for
$U_SW$. There are $K!$ possible $S$ and $2^K K!$ signed permutations,
which proves the bound; different choices can produce the same matrix.
\end{proof}

\paragraph{A rational example.}
Take
\[
 P=\frac15\begin{pmatrix}3&4\\-1&7\end{pmatrix},\qquad
 Q_*=\frac15\begin{pmatrix}3&-4\\4&3\end{pmatrix}.
\]
Then $\det P=1$, $Q_*\in SO(2)$, and
\[
 PQ_*=\begin{pmatrix}1&0\\1&1\end{pmatrix},\qquad
 M_P(Q_*)=\begin{pmatrix}1&0\\2&1\end{pmatrix}.
\]
This is a nonidentity, nondiagonalizable minimizer of the objective.
In contrast,
\[
 M_P(I)=\begin{pmatrix}1&24/25\\14/25&1\end{pmatrix},\qquad
 \rho(M_P(I))=1+\frac{4\sqrt{21}}{25}>1.
\]

\section{Relation to diagonal scaling}

There is a precise link with optimally scaled condition numbers. For fixed
$A\in GL_K(\RR)$, put $s(A)=\rho(|A|\,|A^{-1}|)$ and
$\kappa_\infty(B)=\|B\|_\infty\|B^{-1}\|_\infty$.
Rump \cite[Lemma 2.2 and Theorem 3.1]{rump} proves, without irreducibility
assumptions,
\begin{align*}
 s(A)&=\inf_{D>0\text{ diagonal}}
    \|D^{-1}|A^{-1}|\,|A|D\|_\infty\\
 &=\inf_{D_1,D_2>0\text{ diagonal}}\kappa_\infty(D_1AD_2).
\end{align*}
Here the equality of the spectral radii of $XY$ and $YX$ follows from
$\det(I-zXY)=\det(I-zYX)$, applied to $X=|A|$, $Y=|A^{-1}|$.
It does not require $|A|$ to be invertible.
These are infima: for $T=\left(\begin{smallmatrix}1&0\\1&1\end{smallmatrix}\right)$
and $D=\diag(d_1,d_2)>0$,
\[
 \|D^{-1}|T^{-1}|\,|T|D\|_\infty=1+2d_1/d_2>1,
\]
although the infimum is one. The orthogonal minimum of
Theorem~\ref{thm:minimum}, by contrast, is attained.

\section{The companion row-balancing question}

User3's question 515547 \cite{balance} asks whether, for nonsingular
lower triangular $L$, there are $Q\in O(K)$ and $c>0$ such that
$|Q^\top L^{-1}|e=ce$, where $e=(1,\ldots,1)^\top$.
The answer is affirmative, including without triangularity of $L$.

\begin{corollary}\label{cor:balance}
For every $K\geq1$ and $L\in GL_K(\RR)$, there exist $Q\in SO(K)$ and
$c>0$ satisfying $|Q^\top L^{-1}|e=ce$.
\end{corollary}

\begin{proof}
For $K\geq2$, the Yamabe--Yujobo theorem
\cite[p.~22, Theorem]{yamabe} states that a continuous real-valued
function on the unit sphere in $\RR^K$ takes a common value on some
orthonormal basis. Apply it to
$f(u)=\|u^\top L^{-1}\|_1$. This function is continuous and strictly
positive on the sphere, since $L$ is invertible. For the resulting basis
$q_1,\ldots,q_K$, put $Q=[q_1\ \cdots\ q_K]$ and $c=f(q_i)>0$.
The $i$th component of $|Q^\top L^{-1}|e$ is exactly $f(q_i)$.
Since $f(-u)=f(u)$, a column-sign change can make $\det Q=1$ while
preserving all row sums. For $K=1$, take $Q=1$ and $c=|L^{-1}|$.
\end{proof}

This settles exact balancing and hence the source's fallback question
about approximate regularization. Balancing and spectral minimization
need not use the same $Q$. For $P=L=\diag(1,2)$, every orthogonal
$2\times2$ matrix has
$|Q|=\left(\begin{smallmatrix}a&b\\b&a\end{smallmatrix}\right)$ with
$a,b\geq0$ and $a^2+b^2=1$. Equal inverse-factor row sums require
$a+b/2=b+a/2$, so $a=b=1/\sqrt2$. Then
\[
 M_P(Q)=\begin{pmatrix}1&1/2\\2&1\end{pmatrix},\qquad
 \rho(M_P(Q))=2.
\]
Thus no balancing matrix is a spectral minimizer in this example.

\paragraph{Formal correspondence.}
The main theorem's universal lower bound and existence of an attaining
orthogonal matrix are the statement of \cite[\texttt{result}]{formal},
at the pinned revision listed below. Its definitions use entrywise absolute
values and the complex spectral radius of the real matrix. That source
proves the lower bound by spectral-radius monotonicity and constructs an
attaining matrix by Gram--Schmidt triangularization. The trace bound above
is an alternative elementary proof. The equality classification,
special-orthogonal refinement, examples and balancing corollary have the
mathematical proofs given here.

\paragraph{Availability and license.}
The accompanying source package contains this article's complete LaTeX
source, metadata and a standard-library rational-arithmetic check program,
with build instructions. Original exposition: CC BY 4.0,
\url{https://creativecommons.org/licenses/by/4.0/}.
The MathOverflow questions are attributed to User3 under CC BY-SA 4.0
and are paraphrased here.

\begin{thebibliography}{9}
\bibitem{question}
User3 (MathOverflow user 600125),
\emph{A nontrivial lower bound for the orthogonally transformed absolute product},
MathOverflow question 515549, revision 1, September 25, 2026,
revision GUID \texttt{CB719581-BCE3-4573-931E-15327002F861}.
\url{https://mathoverflow.net/questions/515549}.

\bibitem{formal}
The Omega Institute, \emph{OrthogonalAbsoluteProductMinimum},
Lean source, \texttt{D5/S1/Eigenstructure/}\allowbreak
\texttt{OrthogonalAbsoluteProductMinimum.lean},
theorem \texttt{result}, 2026.
Revision \texttt{fd1633144086085a16f35c8171e2f4b6186a6fbb}.
\href{https://github.com/the-omega-institute/trureturing/blob/fd1633144086085a16f35c8171e2f4b6186a6fbb/D5/S1/Eigenstructure/OrthogonalAbsoluteProductMinimum.lean}{Pinned source}.

\bibitem{qr}
N.~J. Higham, \emph{What Is a QR Factorization?}, November 10, 2020.
\url{https://nhigham.com/2020/11/10/what-is-a-qr-factorization/}.

\bibitem{rump}
S.~M. Rump, \emph{Optimal scaling for $p$-norms and componentwise distance
to singularity}, IMA Journal of Numerical Analysis \textbf{23} (2003), 1--9.
\href{https://doi.org/10.1093/imanum/23.1.1}{doi:10.1093/imanum/23.1.1}.

\bibitem{balance}
User3 (MathOverflow user 600125),
\emph{Can orthogonal transformations balance absolute row sums?},
MathOverflow question 515547, September 25, 2026,
body in revision 1, GUID \texttt{67542193-9900-47F2-A682-20344D440FF0};
revision 2 changes title capitalization.
\url{https://mathoverflow.net/questions/515547}.

\bibitem{yamabe}
Hidehiko Yamabe and Zuiman Yujobo, \emph{On the continuous function defined on a
sphere}, Osaka Mathematical Journal \textbf{2} (1950), 19--22.
\href{https://doi.org/10.18910/4429}{doi:10.18910/4429}.
\end{thebibliography}
\end{document}
